\section{Conclusion}
\label{sec:conclusion}

This thesis evaluated the causal environmental impact of Kenya's settlement schemes, which were implemented to redistribute former colonial land to landless Kenyans.
By employing a spatial regression discontinuity design, the study used high-resolution satellite imagery to measure the impact of settlement schemes on vegetation health and forest cover precisely at the scheme boundaries.

The empirical results show a degradation of a land plot's local ecosystem after its allocation to a settlement scheme.
Within the scheme boundaries, vegetation health is significantly reduced compared to adjacent non-scheme areas.
Contrary to expectations about land tenure security, this thesis could not prove that this negative effect is influenced by the duration of formal land registration.
Furthermore, the assignment is associated with a decrease in forest cover in the surrounding areas, creating a spillover effect.
This effect is strongest just beyond a 300-meter buffer from the settlement boundary, suggesting that the environmental impact of settlement schemes extends beyond the allocated land plots.

This study demonstrates that land redistribution programs create measurable environmental consequences for the local ecosystem.
As those programs are used by many countries to address historical injustices and provide access to land for marginalized populations, the ecological costs should not be underestimated.
Only by mitigating the environmental impact of resettlement can the long-term goal of fighting against externally imposed landlessness be achieved.